\section*{Erd\H{o}s Problem \#910: connected subsets and their number}
\subsection*{Statement and a wording distinction}
The current problem page asks whether every connected subset of Euclidean space contains a nondegenerate connected subset not homeomorphic to itself, and whether a connected subset of \(\mathbb R^n\), \(n\geq2\), must contain more than continuum many connected subsets.
For the first question the original set must also be nondegenerate; a singleton would otherwise be a trivial exception.

There is an important distinction in the second question. Erd\H{o}s's original 1982 discussion says that the connected set itself has dimension greater than one. This is stronger than merely lying in an ambient space of dimension at least two.
Below we settle the literal ambient-dimension version, prove a positive result for sets with interior, and extract a precise cardinality consequence of Rudin's construction. We do not silently transfer these conclusions to the intrinsic-dimension version or to the homeomorphism question.

\subsection*{The literal ambient-dimension version}
Embed \([0,1]\) as a line segment in \(\mathbb R^n\), \(n\geq2\).
Every connected subset of the line is an interval, possibly a point and with either endpoint included or excluded.
To see this, if \(x<z<y\), with \(x,y\) in the set and \(z\) missing, the two intersections with \((-\infty,z)\) and \((z,\infty)\) separate the set.
An interval is determined by its two endpoints and two inclusion choices.
There are at most \(\mathfrak c^2\cdot4=\mathfrak c\) such subsets, and the singletons already give \(\mathfrak c\).
Thus the literal version has a negative answer without additional set-theoretic assumptions.
This example has intrinsic dimension one and does not refute the original higher-dimensional formulation.

\subsection*{A positive result when the set has interior}
If \(A\subseteq\mathbb R^n\), \(n\geq2\), contains an open ball \(B\), then it has exactly \(2^{\mathfrak c}\) connected subsets.
Choose a closed line segment \(L\) strictly inside \(B\), and let \(D=B\setminus L\).
The set \(D\) is path connected: polygonal paths in the ball can be detoured around the segment in a two-dimensional plane; endpoints of the segment lie strictly inside the ball, so the detour can go around an endpoint if necessary.
Also \(D\) is dense in \(B\).
For every subset \(T\subseteq L\),
\[
 D\subseteq D\cup T\subseteq B\subseteq\overline D
\]
implies \(D\cup T\) is connected.
Indeed, a separation of an intermediate set would restrict to a separation of \(D\), and a nonempty side missing \(D\) contradicts density.
The \(2^{\mathfrak c}\) choices of \(T\) give distinct connected subsets of \(A\).
The reverse bound follows from \(|A|\leq\mathfrak c\).
Hence any difficult example for the counting question must have empty interior.

\subsection*{The exact consequence of Rudin's theorem}
The primary theorem of M. E. Rudin (1958), under the continuum hypothesis, constructs a nondegenerate connected plane set \(M\) with this property:
\[
 N\subseteq M,\quad N\text{ connected and nondegenerate}
 \quad\Longrightarrow\quad M\setminus N\text{ is countable}. \tag{R}
\]
We use (R) as an explicit external theorem; its transfinite construction and separation lemmas are not claimed proved here.

It follows that \(M\) has exactly \(\mathfrak c\) connected subsets.
First, \(|M|=\mathfrak c\): choose distinct \(p,q\in M\); the continuous map
\(x\mapsto\|x-p\|\) maps \(M\) onto a connected subset of the line containing \(0\) and \(\|q-p\|\), and hence onto an interval of cardinality \(\mathfrak c\).
The opposite inequality follows from \(M\subseteq\mathbb R^2\).

Every nondegenerate connected subset is determined by a countable complement. The collection of countable subsets of \(M\) has cardinality at most
\[
 |M|^{\aleph_0}
 =\mathfrak c^{\aleph_0}
 =(2^{\aleph_0})^{\aleph_0}
 =\mathfrak c.
\]
Adding the empty set and all singletons does not increase this bound, while the singletons give the matching lower bound.
Notice that the cardinal arithmetic in this deduction does not use CH; CH enters through the construction of \(M\).

Property (R) is a cardinality statement, not a statement that every nondegenerate connected subset is homeomorphic to \(M\).
Such a homeomorphism conclusion would require an additional proof.
Likewise, applying \(M\) to the original intrinsic-dimension version requires the relevant dimension assertion; it has not been established by the argument above.

\subsection*{Attribution and remaining gap}
The sources report known negative resolutions, but this attempt claims only the explicit results proved above and the deduction from (R).
Sources: \url{https://www.erdosproblems.com/910};
Erd\H{o}s, \emph{Some of my favourite problems}, section 4,
\url{https://users.renyi.hu/~p_erdos/1982-33.pdf};
Rudin, \emph{A connected subset of the plane}, Fundamenta Mathematicae 46 (1958), 15--24,
\url{https://matwbn.icm.edu.pl/ksiazki/fm/fm46/fm4612.pdf}.
No novelty is claimed.
The homeomorphism question, the intrinsic-dimension qualification in the original counting question, and a reconstruction of Rudin's transfinite construction remain outside the established argument. Attempt status: unresolved for the complete pair of intended questions.
\subsection*{COMPLETION ESTIMATE}
\textbf{COMPLETION: 45\%}
